\begin{figure}[htbp]
		\centering
		\subfloat[][\emph{sistema di partenza}]{%
			\begin{tikzpicture}[>=Stealth, line width=0.8pt, scale=0.80]
				\fill (0,0) circle (1.5pt) node[below] {$Q$};
				\fill (1.0,2.0) circle (1.5pt) node[above right] {$P$};
				% vettore QP
				\draw[->, thin, gray!60] (0,0) -- (1.0,2.0)
				node[midway, left] {$\overrightarrow{QP}$};
				% forza in P
				\draw[->, thick] (1.0,2.0) -- (2.8,2.0)
				node[above] {$\fb$};
		\end{tikzpicture}}
		\qquad
		\subfloat[][\emph{aggiunta della coppia equilibrata in $Q$}]{%
			\begin{tikzpicture}[>=Stealth, line width=0.8pt, scale=0.80]
				\fill (0,0) circle (1.5pt) node[below] {$Q$};
				\fill (1.0,2.0) circle (1.5pt) node[above right] {$P$};
				% vettore QP
				\draw[->, thin, gray!60] (0,0) -- (1.0,2.0)
				node[midway, left] {$\overrightarrow{QP}$};
				% forza originale in P
				\draw[->, thick] (1.0,2.0) -- (2.8,2.0)
				node[above] {$\fb$};
				% forza +f in Q
				\draw[->, thick] (0,0) -- (1.8,0)
				node[below] {$\fb$};
				% forza -f in Q
				\draw[->, thick] (0,0) -- (-1.8,0)
				node[below] {$-\fb$};
				% braccio b
				\draw[<->, thin] (2.9,0) -- (2.9,2.0)
				node[midway, right] {$b$};
		\end{tikzpicture}}
		\qquad
		\subfloat[][\emph{sistema equivalente}]{%
			\begin{tikzpicture}[>=Stealth, line width=0.8pt, scale=0.80]
				\fill (0,0) circle (1.5pt) node[below] {$Q$};
				\fill (1.0,2.0) circle (1.5pt) node[above right] {$P$};
				% forza risultante in Q
				\draw[->, thick] (0,0) -- (1.8,0)
				node[below] {$\fb$};
				% coppia correttiva (freccia semicircolare oraria)
				\draw[->, thick] (-0.6,0) arc[start angle=180,
				end angle=90, radius=0.6]
				node[above left] {$\mu(Q)$};
				% braccio b
				\draw[<->, thin] (2.9,0) -- (2.9,2.0)
				node[midway, right] {$b$};
		\end{tikzpicture}}
		\caption{Trasporto di una forza in un punto non appartenente
			alla linea d'azione.
			\textbf{\textit{(a)}}~Sistema originale: forza $\fb$
			applicata in $P$.
			\textbf{\textit{(b)}}~Aggiunta in $Q$ della coppia
			equilibrata $(\fb,\,-\fb)$: la forza $(P,\,\fb)$ e
			$(Q,\,-\fb)$ formano una coppia di momento
			$\bm{\upmu} = \protect\overrightarrow{QP}\times\fb$,
			con braccio $b$.
			\textbf{\textit{(c)}}~Sistema equivalente: forza $\fb$
			applicata in $Q$ più la coppia correttiva oraria di intensità 
			$\mu(Q) = |\fb|\,b$.}
		\label{fig:trasporto_forza}
	\end{figure}
